\documentclass[a4paper]{article}
% Español (México) con XeLaTeX: fontspec para fuentes Unicode, polyglossia
% para la división silábica y los nombres del documento en la variante mexicana.
\usepackage{fontspec}
\usepackage{polyglossia}
\setdefaultlanguage[variant=mexican]{spanish}
% Los nombres chinos van como «pinyin (original)»: xeCJK da a los caracteres
% Han su propia fuente; sin ella XeLaTeX los omite con un aviso.
% FandolSong no cubre algunos caracteres de nombres (U+73FA); WenQuanYi Zen Hei sí.
\usepackage[AutoFallBack=true]{xeCJK}
\setCJKmainfont{FandolSong-Regular.otf}
\setCJKfallbackfamilyfont{\CJKrmdefault}{WenQuanYi Zen Hei}
% polyglossia no traduce los nombres de listings.
\AtBeginDocument{\providecommand{\lstlistingname}{}\renewcommand{\lstlistingname}{Listado}%
  \providecommand{\lstlistlistingname}{}\renewcommand{\lstlistlistingname}{Índice de listados}}
\usepackage{amsmath, amssymb}
\usepackage{graphicx}
\usepackage[margin=2.35cm]{geometry}
\usepackage[most]{tcolorbox}
\usepackage{booktabs}
\usepackage{tabularx}
\usepackage{longtable}
\usepackage{array}
\usepackage{float}
\usepackage{xcolor}
\usepackage{hyperref}
\usepackage{fancyhdr}
\hypersetup{colorlinks=true, linkcolor=blue!55!black, urlcolor=blue!60!black}
\graphicspath{{./}{figures/}}
\setlength{\parskip}{0.55em}
\setlength{\parindent}{2em}
\setlength{\emergencystretch}{3em}
\renewcommand{\arraystretch}{1.18}
\newtcolorbox{knowledgebox}[1]{enhanced, colback=blue!5!white, colframe=blue!70!black, colbacktitle=blue!70!black, coltitle=white, fonttitle=\bfseries, title={#1}, sharp corners, breakable}
\newtcolorbox{importantbox}[1]{enhanced, colback=yellow!10!white, colframe=yellow!75!black, colbacktitle=yellow!75!black, coltitle=black, fonttitle=\bfseries, title={#1}, sharp corners, breakable}
\newtcolorbox{warningbox}[1]{enhanced, colback=red!5!white, colframe=red!70!black, colbacktitle=red!70!black, coltitle=white, fonttitle=\bfseries, title={#1}, sharp corners, breakable}
\pagestyle{fancy}
\fancyhf{}
\fancyfoot[C]{\thepage}
\fancyfoot[R]{\footnotesize\color{gray}\href{https://github.com/hqhq1025/ai-course-notes}{AI Course Notes}}
\renewcommand{\headrulewidth}{0pt}
\renewcommand{\footrulewidth}{0pt}
\newcommand{\videourl}{https://www.youtube.com/watch?v=x8qdqWIVVTA}
\newcommand{\repourl}{https://github.com/hqhq1025/ai-course-notes}
\begin{document}
\begin{titlepage}
\centering
{\Large Notas de entrevista técnica\par}
\vspace{1.0cm}
{\huge\bfseries Elys, dobles cibernéticos y redes sociales de IA: cómo fluye el Context\par}
\vspace{0.6cm}
{\Large Zhang Xiaojun conversa con Tristan (fundador de Ziran Xuanze (自然选择))\par}
\vspace{0.4cm}
{\large Elaborado a partir del video público del Zhang Xiaojun Podcast, una transcripción generada en GPU y diagramas conceptuales propios\par}
\vspace{0.3cm}
{\large 2026-04-10\par}
\vspace{0.9cm}
\includegraphics[width=0.88\textwidth,height=0.42\textheight,keepaspectratio]{cover.jpg}\par
\vfill
\begin{tcolorbox}[width=0.92\textwidth, colback=black!2!white, colframe=black!55, sharp corners]
\textbf{Título del video}: 135. Conversación con Tristan, fundador de Ziran Xuanze: Elys, dobles cibernéticos, el alma, la obtención y el flujo del Context, y las redes sociales de IA\par
\textbf{Canal del video}: Zhang Xiaojun Podcast\par
\textbf{Duración del video}: 01:52:55\par
\textbf{Enlace del video}: \href{\videourl}{\nolinkurl{\videourl}}\par
\textbf{Notas sobre el material}: YouTube no ofrece subtítulos disponibles; el texto se elaboró a partir de una transcripción con faster-whisper large-v3 en CUDA, los capítulos del video, la portada y diagramas didácticos propios. Las tomas fijas de la entrevista no se repiten en el cuerpo del texto.\par
\textbf{Página del proyecto}: \href{\repourl}{\nolinkurl{\repourl}}\par
\end{tcolorbox}
\end{titlepage}
\tableofcontents
\newpage
\section{Introducción: la socialización con IA no es «una IA que platica contigo»}

Esta entrevista gira en torno a dos productos de Ziran Xuanze (自然选择): el juego con IA Eve y la red social con IA Elys. La narrativa central que plantea Tristan es que «los seres humanos están demasiado solos», pero no se trata de un simple producto de acompañamiento emocional. La ambición mayor de Elys es que el doble de IA salga activamente al mundo y ayude al usuario a convertir su propio context en conexiones reales. Intenta responder una pregunta que dejó pendiente la era de internet: por qué las redes sociales siempre comprimen a personas complejas en etiquetas de baja dimensión, y si la IA puede, por primera vez, hacer que fluya el context de alta dimensión.

\begin{importantbox}{Tesis central de este episodio}
La socialización tradicional en internet depende de etiquetas de baja dimensión y de las acciones que el propio usuario emprende; las redes sociales con IA intentan que el usuario construya su subjetividad, que su doble cibernético, con su context a cuestas, se exprese, recomiende y rompa el hielo por iniciativa propia, y que al final conduzca de vuelta a conexiones entre personas reales.
\end{importantbox}

\begin{knowledgebox}{Nota sobre la estrategia visual}
Este video es una entrevista con encuadre fijo, sin slides didácticas, pizarrón ni demostraciones del producto. Conforme al estándar de este repositorio para pódcasts, el cuerpo del texto no repite fotogramas de las personas; la portada sirve para identificar la fuente, y el cuerpo del texto presenta la lógica del producto mediante diagramas conceptuales y tablas.
\end{knowledgebox}

\subsection{Resumen de la sección}

Elys no es un simple Tinder con IA, ni tampoco un juguete en el que una IA platica con otra IA. Sus variables clave son el context, la subjetividad, el intercambio de privacidad, la tasa de conexión entre personas reales y la forma de producto proactive.
\section{El flujo del context: el mecanismo de fondo de Elys}

Tristan insiste varias veces en que el núcleo de Elys es el flujo del context. El context no son las etiquetas de intereses del internet tradicional, sino las experiencias de una persona, su forma de expresarse, sus valores, su sentido estético, sus relaciones, sus obras pasadas y su intención actual. Elys busca que el usuario entregue ese context a su doble digital, y que el doble salga por iniciativa propia al entorno social a expresarse, comentar, recomendar y romper el hielo, para que al final traiga de vuelta conexiones entre personas reales.

\begin{figure}[H]
\centering
\includegraphics[width=0.94\textwidth]{context-flow.png}
\caption{El ciclo cerrado del flujo del context en Elys.}
\end{figure}

\begin{knowledgebox}{Lectura de la figura: el context no es una página de datos estática}
En la figura, el context circula entre el usuario, el doble cibernético, las acciones de la IA, las conexiones entre personas reales y la retroalimentación. No es una etiqueta de baja dimensión como «me gusta la música», sino un estado de alta dimensión que permite que el doble se parezca a ti, hable en tu nombre e identifique por ti oportunidades de relación.
\end{knowledgebox}

\subsection{Por qué el internet tradicional se basa en etiquetas de baja dimensión}

Las redes sociales del internet tradicional suelen depender del apodo, la foto de perfil, la región, la edad, los intereses, las relaciones de seguimiento y los registros de comportamiento. Los sistemas de recomendación pueden distribuir contenido con base en estas etiquetas, pero difícilmente entienden los valores, el sentido del humor, la forma de expresarse, la experiencia de vida y los objetivos de relación de una persona. El audio largo, los textos extensos, los historiales de chat, las obras y la retroalimentación social real apenas en la era de la IA empiezan a tener la posibilidad de entenderse de manera unificada.

\begin{importantbox}{El context es el activo de las redes sociales con IA}
La inteligencia puede volverse gradualmente accesible para todos por igual, pero el context no. Quien posea un context de usuario de más largo plazo, más auténtico y cuyo uso pueda autorizarse, estará más cerca del activo central de una red social con IA.
\end{importantbox}

\subsection{El volante de inercia del context}

El usuario aporta un context inicial y el doble empieza a actuar; las acciones generan interacciones, y las interacciones producen nueva retroalimentación; esa retroalimentación, a su vez, hace que el doble se parezca más al usuario. Lo más difícil de este volante de inercia es el arranque en frío: al principio el usuario no confía en el producto y no quiere entregar su privacidad; el producto tiene que mostrarle los beneficios reales que aporta el context antes de que lo abandone.

\begin{warningbox}{Los puntos frágiles del volante de inercia del context}
Si el usuario entrega muy poco context, el doble no se parece a él y el producto no le dice nada; si el producto pide demasiado, el usuario se preocupa por su privacidad y por ser representado de forma equivocada. La dificultad de Elys consiste justamente en encontrar un punto de arranque entre los beneficios y los riesgos.
\end{warningbox}

\subsection{Resumen de la sección}

La base de Elys no es el chat, sino el ciclo del context. Que las redes sociales con IA funcionen depende de que el context pueda obtenerse de forma segura, fluir de manera eficaz y convertirse en conexiones entre personas reales.
\section{Subjetividad: el significado de ingeniería de «ser uno mismo» en la era de la IA}

Uno de los conceptos más fuertes de la entrevista es «construir la subjetividad». Tristan dice que, en la era de la IA, lo único importante quizá sea «ser uno mismo», pero aquí «ser uno mismo» no es un mensaje de autoayuda, sino lograr que la IA sepa quién eres: tus valores, tu estética, tus activos históricos, tu forma de expresarte, tus objetivos actuales y tus límites. Cuanto más clara sea la subjetividad, mejor podrá el avatar actuar en tu nombre.

\begin{figure}[H]
\centering
\includegraphics[width=0.94\textwidth]{subjectivity-stack.png}
\caption{Construcción de la subjetividad: la pila de contexto que permite a la IA actuar en tu nombre.}
\end{figure}

\begin{knowledgebox}{Lectura de la figura: por qué la subjetividad es un problema de ingeniería}
En la figura, la subjetividad se compone de los activos históricos, los valores y la estética, los objetivos actuales, los límites de conducta y la expresión agentic. No es una personalidad abstracta, sino un conjunto de entradas que determina si el avatar de IA se parece a ti, si puede tomar decisiones por ti y si sabe cuándo detenerse.
\end{knowledgebox}

\subsection{Por qué «ser uno mismo» cobra importancia}

Cuando la IA puede escribir artículos, publicar comentarios, reconocer personas, generar contenido y recomendar conexiones por iniciativa propia, lo único insustituible del usuario puede pasar a ser «quién es». Si una persona ha acumulado durante mucho tiempo obras, opiniones, estética y relaciones, la IA puede amplificarla con mayor precisión; si una persona carece de subjetividad, la IA solo puede generar una expresión promediada.

\begin{importantbox}{La subjetividad es el activo personal del futuro}
En la era de internet, el tráfico y los seguidores eran activos; en la era de las redes sociales con IA, la subjetividad que un modelo puede comprender y representar se convertirá en un activo. Incluye obras, expresión, valores, relaciones y límites.
\end{importantbox}

\subsection{Los límites del avatar}

Cuanto más se parece el avatar al usuario, mayor es el riesgo: puede decir algo equivocado, expresarse en nombre del usuario sin autorización o filtrar información privada. Por eso la subjetividad debe incluir límites: qué contenido puede hacerse público, qué contenido requiere confirmación y qué contenido el avatar nunca puede usar.

\begin{warningbox}{Un avatar sin límites no es un producto, es un riesgo}
Si un avatar solo busca «parecerse a ti», pero carece de límites de permisos, de auditoría y de mecanismos de revocación, se convierte en un agente que el usuario no puede controlar. Los productos sociales con IA deben construir a la vez la semejanza y la sensación de control.
\end{warningbox}

\subsection{Resumen de la sección}

La subjetividad es la entrada central de las redes sociales con IA. Permite que el avatar actúe en nombre del usuario y, al mismo tiempo, exige que el producto ofrezca límites y mecanismos de control.
\section{El intercambio de privacidad: ceder context a cambio de conexiones}

El problema de producto de Elys es lograr que los usuarios estén dispuestos a ceder su context. El conductor señala que los usuarios adoptan de forma natural una actitud defensiva: no quieren que una organización comercial sepa demasiado de ellos, no quieren que su doble digital diga cosas indebidas y tampoco tienen claro qué beneficio obtendrán al ceder su context. La respuesta de Tristan es que el producto debe hacer que el usuario perciba los beneficios del ciclo virtuoso del context antes de que lo abandone.

\begin{figure}[H]
\centering
\includegraphics[width=0.95\textwidth]{context-privacy-tradeoff.png}
\caption{El intercambio entre context y privacidad: beneficios y riesgos de ceder poco, una cantidad moderada o demasiado.}
\end{figure}

\begin{knowledgebox}{Lectura de la figura: más privacidad no siempre es mejor, y más context tampoco}
Ceder poco context implica poco riesgo, pero el usuario no percibe el valor del producto; ceder una cantidad moderada permite mejores emparejamientos y conexiones; ceder demasiado puede hacer que la situación se salga de control. Un buen producto debe permitir que el usuario vea los beneficios y que controle el alcance de la autorización.
\end{knowledgebox}

\subsection{Almacenamiento local y confianza}

En la entrevista se menciona que los usuarios probablemente siempre confiarán más en que su información privada esencial se guarde localmente. Esto significa que los productos sociales de IA deben diseñar memoria local, autorizaciones revocables, registros transparentes, niveles de privacidad y un comportamiento del doble digital que pueda explicarse. De lo contrario, cuanto más esencial sea el context, menos se atreverá el usuario a cederlo.

\subsection{Conexiones reales frente al bullicio de la IA}

Tristan afirma con claridad que tiene poco sentido que las IA publiquen y comenten entre sí. Lo que le importa al usuario es si detrás hay una persona real. La IA puede abrir camino y reducir la barrera para romper el hielo, pero la estrella polar debe ser la tasa de conexiones entre personas reales y la proporción de acciones realizadas por personas reales.

\begin{importantbox}{Métrica estrella polar: tasa de conexiones entre personas reales}
Una red social de IA no existe para que las IA se entretengan entre ellas, sino para que las personas reales se conecten mejor. Por eso, la estrella polar de Elys no es el DAU ni el número de mensajes de IA, sino la tasa de conexiones entre personas reales, la proporción de acciones realizadas por personas reales y si los usuarios encuentran a personas que de verdad valen la pena.
\end{importantbox}

\subsection{Resumen de la sección}

Las redes sociales de IA deben resolver el intercambio entre context y privacidad. Sin context, el doble digital no sirve; sin confianza, el usuario no cede su context; sin conexiones entre personas reales, las redes sociales de IA pierden su sentido.
\section{Red social de IA frente a redes sociales de internet}

A Elys le preguntaron si es una versión de Tinder con IA. La respuesta de Tristan es que no. El núcleo de Tinder es el emparejamiento de baja dimensión y la selección deslizando; lo que Elys quiere construir es un emparejamiento de context de alta dimensión y una socialización proactiva a través del doble digital. En las redes sociales tradicionales de internet, la persona tiene que expresarse, filtrar y romper el hielo por iniciativa propia; la socialización con IA busca que el doble digital se exprese y explore primero en tu lugar, y después te traiga de vuelta las posibles conexiones reales.

\begin{figure}[H]
\centering
\includegraphics[width=0.92\textwidth]{ai-social-vs-internet-social.png}
\caption{Diferencias entre las redes sociales de internet y la red social de IA.}
\end{figure}

\begin{knowledgebox}{Lectura de la figura: la diferencia de la socialización con IA no es solo «agregar IA»}
A la izquierda, las redes sociales de internet dependen de etiquetas de baja dimensión y del comportamiento proactivo del usuario; a la derecha, la socialización con IA depende de un context de alta dimensión y de la acción proactiva del doble digital. La diferencia real es que la IA reduce el costo de conectar, no que la interfaz se convierta en una ventana de chat.
\end{knowledgebox}

\subsection{Por qué no es un Tinder con IA}

Tinder se parece más a una criba rápida de candidatos para una relación, centrada en los datos visibles del perfil y en la elección inmediata. Elys quiere construir más bien una red de context: los valores, las aficiones y las formas de expresarse de una persona, por sutiles que sean o aunque no convenga decirlos en público, también pueden encontrar, mediante la IA, a alguien con quien resuenen. Lo que persigue es una conexión con menos fricción y de mayor dimensión.

\begin{warningbox}{Un error peligroso de la socialización con IA}
Si la socialización con IA solo genera una animación falsa, desgastará la confianza con rapidez. El usuario puede aceptar que la IA lo ayude a romper el hielo, pero el valor real sigue proviniendo de la conexión entre personas reales.
\end{warningbox}

\subsection{Quién posee más context}

En la entrevista se menciona que ChatGPT, Meta, Doubao (豆包) y otros podrían poseer una gran cantidad de context. El juicio de Tristan es que la inteligencia se irá distribuyendo de forma equitativa, pero el context no. Las empresas de modelos tienen la capacidad de los modelos, las plataformas sociales tienen las redes de relaciones, ChatGPT tiene la memoria personal y el historial de conversaciones, y las empresas emergentes de socialización con IA intentan construir productos especializados en torno al ciclo virtuoso del context.

\begin{figure}[H]
\centering
\includegraphics[width=0.92\textwidth]{context-competition.png}
\caption{Variables de competencia en la socialización con IA: capacidad de los modelos, redes de relaciones sociales, activos de Context y agencia del usuario.}
\end{figure}

\begin{knowledgebox}{Lectura de la figura: por qué el Context es más escaso que el modelo}
Los cuatro tipos de actores de la figura tienen cada uno su ventaja: las empresas de modelos tienen la inteligencia, las plataformas sociales tienen las relaciones, las empresas emergentes tienen un nuevo paradigma y el propio usuario tiene la agencia. Si la inteligencia se convierte poco a poco en una mercancía, el context de largo plazo y las relaciones entre personas reales pasarán a ser los activos más escasos.
\end{knowledgebox}

\subsection{Resumen de la sección}

El núcleo de la red social de IA no es una nueva generación de algoritmos de emparejamiento, sino el paso de una socialización basada en etiquetas de baja dimensión a una socialización basada en context de alta dimensión. Su objetivo es que la IA reduzca el costo de conectar a personas reales, no sustituir esa conexión.
\section{Proactive: la diferencia esencial entre los productos de IA y los productos de internet}

Tristan considera que el mayor cambio de interacción en la era de la IA no es LUI ni GUI, sino lo proactive: por fin hay algo que puede actuar por iniciativa propia en tu nombre. OpenClaw permitió experimentar lo proactive en la capa de herramientas; Elys, en cambio, lleva lo proactive a la capa social y hace que el doble digital salga por iniciativa propia a relacionarse con el mundo.

\begin{figure}[H]
\centering
\includegraphics[width=0.92\textwidth]{proactive-product-shift.png}
\caption{Paradigmas de producto: de Reactive a Assistive y de ahí a Proactive.}
\end{figure}

\begin{knowledgebox}{Lectura de la figura: Proactive cambia el modelo mental del producto}
Los productos Reactive requieren que el usuario tome la iniciativa; los productos Assistive ayudan mientras el usuario opera; los productos Proactive pueden explorar oportunidades en nombre del usuario cuando este no hace nada por iniciativa propia. El potencial de Elys proviene del tercer nivel: el doble digital socializa por iniciativa propia y le trae conexiones al usuario.
\end{knowledgebox}

\subsection{La competencia con las empresas de modelos}

Tristan considera que la IA social no necesariamente requiere entrenar un modelo base propio. La razón es que el núcleo de una red social no es solo la inteligencia, sino el context, el paradigma de producto y la conexión entre personas reales. Las empresas de modelos se enfocan en la inteligencia en sí y no necesariamente tienen lo social como primera prioridad; si una empresa de aplicaciones logra controlar el volante de inercia del context, también puede tener una oportunidad independiente.

\begin{importantbox}{¿Dónde está la barrera de entrada de una empresa de aplicaciones?}
No se trata de «yo también puedo llamar al modelo», sino de si tengo un context propio, relaciones reales, un modelo mental de producto, retroalimentación de datos y la confianza de los usuarios. Sin esto, las empresas de modelos absorberán la aplicación; con esto, la aplicación puede generar sus propios efectos de red.
\end{importantbox}

\subsection{Los seres humanos están demasiado solos}

Al final, la entrevista vuelve al núcleo espiritual: los seres humanos están demasiado solos. Eve intenta resolver la soledad mediante la compañía de la IA; Elys intenta resolverla mediante la IA como promotora de conexiones entre personas reales. Lo primero es una relación entre la IA y la persona; lo segundo es la IA ayudando a que las personas establezcan relaciones entre sí. Esta diferencia determina que Elys no sea un producto de compañía, sino un producto de red social.

\begin{warningbox}{El valor emocional no es una necesidad menor}
Socializar, tener compañía, sentirse comprendido y conectar son necesidades humanas frecuentes y profundas. Si un producto de IA logra atender mejor estas necesidades de forma segura, auténtica y controlable, puede estar más cerca del usuario común que un producto puramente utilitario. Pero el riesgo también es mayor, porque toca la privacidad, la identidad y las relaciones.\end{warningbox}

\subsection{Resumen de la sección}

Lo proactive es lo que distingue a los productos de IA de los productos de internet. El núcleo de Elys no es que la IA converse para acompañar, sino que el doble digital de IA reduzca por iniciativa propia la fricción social en nombre del usuario y, en última instancia, le traiga conexiones con personas reales.

\section{Dificultades del diseño del producto: arranque en frío, doble arranque en frío y proporción de personas reales}

Elys no es un producto que funcione por sí solo con que el modelo sea lo bastante capaz. Debe resolver dos arranques en frío al mismo tiempo: el primero es el arranque en frío de usuarios, es decir, cómo lograr que suficientes personas entren en la red; el segundo es el arranque en frío de context, es decir, cómo lograr que los usuarios estén dispuestos a aportar información sobre sí mismos en cantidad suficiente, auténtica y utilizable. Solo si los usuarios y el context arrancan a la vez puede el doble digital generar conexiones reales.

\subsection{Doble arranque en frío}

\noindent\begin{longtable}{p{0.22\textwidth}p{0.34\textwidth}p{0.35\textwidth}}
\toprule
Tipo de arranque en frío & Problema central & Soluciones posibles \\
\midrule
Arranque en frío de usuarios & Sin suficientes personas reales, la interacción entre dobles parece girar en vacío & Arranque en círculos pequeños, acceso por invitación, comunidades semilla, difusión basada en escenarios de uso. \\
Arranque en frío de context & Los usuarios no quieren o no saben entregar context de alta calidad & Onboarding de baja fricción, beneficio inmediato, autorización revocable, guía con ejemplos. \\
Arranque en frío de confianza & A los usuarios les preocupa que la empresa controle su privacidad o que el doble los represente mal & Almacenamiento local, permisos por niveles, auditoría de acciones, confirmación de acciones sensibles. \\
Arranque en frío de la red & Hay personas pero no relaciones; hay contenido pero no conexiones & La IA rompe primero el hielo y después guía a las personas reales a confirmar e interactuar. \\
\bottomrule
\end{longtable}

\begin{importantbox}{La esencia del doble arranque en frío}
Los productos sociales tradicionales arrancan en frío sobre todo usuarios y contenido; la IA social además tiene que arrancar en frío el context. Sin context, la IA no se parece a ti; sin personas reales, la actividad de la IA no tiene sentido.
\end{importantbox}

\subsection{Por qué importa la proporción de personas reales}

La métrica estrella polar de Elys es la tasa de conexiones entre personas reales, o la proporción de acciones realizadas por personas reales. Este punto es decisivo: la IA puede reducir la barrera para interactuar, pero no puede hacer que en la red solo quede IA. Si cada vez que el usuario abre la aplicación solo ve respuestas de IA, se decepcionará, porque el sentido de lo social sigue proviniendo de las personas reales. El papel de la IA debe ser amplificar las conexiones entre personas reales, no sustituirlas.

\begin{warningbox}{La «trampa del bullicio» en la IA social}
Si el producto cuenta como crecimiento los comentarios, los «me gusta» y las publicaciones hechos por IA, los datos pueden verse bien a corto plazo, pero los usuarios descubrirán pronto que no hay personas reales. Lo que más teme una red social es la prosperidad ficticia, porque destruye la confianza.
\end{warningbox}

\subsection{Resumen de la sección}

La dificultad del producto Elys no está en «si puede generar conversaciones», sino en si puede poner en marcha al mismo tiempo a los usuarios, el context, la confianza y la red de personas reales. La tasa de conexiones entre personas reales es una métrica más cercana a lo esencial que el número de mensajes.
\section{Modelo de privacidad y permisos: cómo debe restringirse al doble digital}

Para que un doble digital salga al mundo de forma proactiva en nombre del usuario, necesita un modelo de permisos. Qué context puede consultar un doble, con quién puede hablar, qué puede expresar en representación del usuario, qué acciones puede completar de forma automática y cuáles requieren confirmación: todas son cuestiones de fondo del producto. Sin un modelo de permisos, cuanto más proactivo sea el doble, más peligroso resulta.

\subsection{Cuatro niveles de permisos}

\noindent\begin{tabularx}{\textwidth}{p{0.22\textwidth}p{0.34\textwidth}X}
\toprule
Nivel de permisos & Qué puede hacer el doble & Control de riesgos \\
\midrule
Nivel de lectura & Leer los datos del usuario, sus expresiones previas, sus preferencias y sus relaciones & El usuario puede consultar, eliminar y autorizar por niveles. \\
Nivel de borrador & Generar comentarios, mensajes privados y justificaciones de recomendaciones, pero sin enviarlos & Se publica tras la confirmación del usuario. \\
Nivel de acciones de bajo riesgo & Dar «me gusta», guardar en favoritos, interacciones públicas de bajo riesgo & Límites de frecuencia y registros de auditoría. \\
Nivel de acciones de alto riesgo & Conversaciones privadas, opiniones sensibles, compromisos en las relaciones, acciones comerciales & Requiere confirmación humana obligatoria y debe poder revertirse. \\
\bottomrule
\end{tabularx}

\begin{knowledgebox}{El modelo de permisos también es experiencia de producto}
La condición para que el usuario esté dispuesto a entregar su context es saber que el doble no se saldrá de control. Unos permisos claros no limitan el producto: hacen que el usuario se atreva a usarlo.
\end{knowledgebox}

\subsection{La propiedad del context}

La cuestión más delicada de las redes sociales con IA es a quién pertenece realmente el context. Las conversaciones, las obras, las preferencias sociales y los datos de relaciones del usuario no pueden ser solo un activo de la plataforma. A largo plazo, el usuario necesita derechos sobre su context para poder trasladarlo, exportarlo, eliminarlo y guardarlo localmente. De lo contrario, cuanto mejor conozca la plataforma al usuario, más inquieto se sentirá este.

\begin{importantbox}{Lista de derechos sobre el context}
El usuario debería tener, como mínimo, derecho de consulta, derecho a una explicación, derecho de eliminación, derecho de exportación, derecho de autorización por niveles y derecho a confirmar las acciones sensibles. Si una red social con IA quiere construir confianza a largo plazo, debe convertir estos derechos en capacidades del producto.
\end{importantbox}

\subsection{Resumen de la sección}

Cuanto más proactivo es el doble digital, más importante se vuelve el modelo de permisos. La base de la confianza en las redes sociales con IA no es que el modelo hable mejor, sino que el usuario pueda controlar su propio context y el comportamiento de su doble.
\section{La lógica competitiva de Elys frente a ChatGPT, Meta y Doubao (豆包)}

En la entrevista aparece un juicio importante: la capacidad de los modelos tenderá a igualarse, pero el context no. Las empresas de modelos tienen una inteligencia potente y los puntos de entrada; las plataformas sociales tienen las redes de relaciones y el comportamiento en torno al contenido; las empresas emergentes como Elys tienen un nuevo paradigma y un producto enfocado. Quién gane no depende solo de quién tenga el modelo más potente, sino de quién logre acumular el context más útil, más confiable y más capaz de convertirse en conexiones entre personas reales.

\subsection{Las trayectorias de cuatro tipos de competidores}

\noindent\begin{longtable}{p{0.20\textwidth}p{0.35\textwidth}p{0.36\textwidth}}
\toprule
Competidor & Ventajas & Debilidades \\
\midrule
ChatGPT / OpenAI & Punto de entrada conversacional, memoria personal, modelo potente & Lo social no es su línea principal de negocio; su posicionamiento en la mente del usuario como producto de entretenimiento y afectivo no necesariamente es fuerte. \\
Meta / plataformas sociales & Redes de relaciones, contenido, identidad y distribución & Puede resultar difícil volver AI-native las estructuras sociales heredadas. \\
Doubao / grandes empresas tecnológicas chinas & Punto de entrada al consumidor final, ecosistema de contenido, escenarios locales & Sigue siendo incierto si consideran el context un eje estratégico. \\
Elys / empresas emergentes & Diseño desde cero con el paradigma de la red social con IA, enfoque en la conexión entre personas reales & Fuerte presión de arranque en frío en usuarios, confianza y context. \\
\bottomrule
\end{longtable}

\begin{warningbox}{Las empresas de modelos no ganan necesariamente por naturaleza}
Si el núcleo del producto es la respuesta inteligente a preguntas, la ventaja de las empresas de modelos es enorme; si el núcleo del producto son las relaciones entre personas reales, el volante de inercia del context y el posicionamiento social en la mente del usuario, las empresas de aplicaciones todavía pueden tener una ventana de oportunidad. Lo decisivo es si el producto se mantiene lejos de la línea principal de negocio de las empresas de modelos.
\end{warningbox}

\subsection{Por qué OpenClaw y Elys pertenecen a la misma narrativa}

OpenClaw hace sentir al usuario que «hay algo que trabaja por mí por iniciativa propia»; Elys hace sentir al usuario que «hay un doble que enfrenta el mundo por mí por iniciativa propia». Ambos encarnan el posicionamiento de producto del proactive agent. La diferencia es que OpenClaw se dirige a las herramientas y a la computadora local, mientras que Elys se dirige a las personas y a las relaciones sociales.

\begin{importantbox}{La misma narrativa: la proactividad}
El cambio de fondo en los productos de IA no es que la interfaz pase de GUI a LUI, sino que pasan de reactive a proactive. Los modelos empiezan a explorar, ejecutar y traer resultados por iniciativa propia, y con ello los productos pasan de ser herramientas a ser agentes que actúan en nombre del usuario.
\end{importantbox}

\subsection{Resumen de la sección}

La competencia en la red social con IA no es solo una competencia de capacidad de los modelos. El context, las redes de relaciones, los permisos, el atractivo del producto y la tasa de conexión entre personas reales determinarán, todos ellos, el panorama final.
\section{Historia de la evolución de las redes sociales: de las plataformas de contenido a la socialización AI Native}

La trayectoria de Tristan va de una plataforma de contenido de audio, la socialización entre desconocidos, los juegos de citas y la compañía con IA, hasta Elys. Este recorrido puede verse como una línea de evolución de los productos sociales: primero, la distribución de contenido; después, el emparejamiento entre personas; luego, personajes virtuales que aportan valor emocional, y por último, dobles digitales de IA que ayudan de forma proactiva a conectar a personas reales.

\subsection{Trayectoria de evolución}

\noindent\begin{longtable}{p{0.22\textwidth}p{0.34\textwidth}p{0.35\textwidth}}
\toprule
Etapa & Mecanismo central & Limitaciones \\
\midrule
Plataforma de contenido & Distribución personalizada de audio/contenido & Conecta a las personas con el contenido, no a las personas entre sí. \\
Socialización entre desconocidos & Etiquetas de baja dimensión y emparejamiento inmediato & Romper el hielo es costoso; la calidad de las relaciones es inestable. \\
Juegos de citas/compañía con IA & El contenido o la IA aporta valor emocional & El interlocutor puede no ser una persona real. \\
Red social de IA & El doble digital, con su context, conecta de forma proactiva a personas reales & Privacidad, confianza, proporción de personas reales y dificultad del arranque en frío. \\
\bottomrule
\end{longtable}

\begin{knowledgebox}{Por qué Elys no surgió de la nada}
Elys hereda la experiencia de la recomendación de contenido, la socialización entre desconocidos, los juegos de citas y la compañía con IA, pero intenta llevar el objetivo de «consumir contenido/compañía» a «conectar a personas reales».
\end{knowledgebox}

\subsection{Resumen de la sección}

La socialización AI Native no aparece de la nada: es la confluencia de tres líneas, las plataformas de contenido, el emparejamiento social y la compañía con IA. La novedad de Elys consiste en que el doble digital convierte de forma proactiva el context en conexiones entre personas reales.
\section{Riesgos del producto: un mundo de bajo daño también puede convertirse en uno de alto daño}

La visión de Elys es un mundo de bajo daño: menos fricción y conexiones más acertadas. Sin embargo, cualquier producto social puede causar daño, y la socialización mediada por IA amplifica ese riesgo. El doble digital puede representar mal al usuario, el context puede filtrarse, la IA puede fabricar una actividad social falsa y el algoritmo puede empujar al usuario hacia una exposición excesiva o hacia relaciones poco sanas.

\subsection{Lista de riesgos}

\noindent\begin{tabularx}{\textwidth}{p{0.22\textwidth}p{0.34\textwidth}X}
\toprule
Riesgo & Manifestación & Mecanismo necesario \\
\midrule
Representación errónea & El doble digital dice cosas que el usuario no respalda & Vista previa, posibilidad de retractarse y confirmación de acciones sensibles. \\
Filtración de privacidad & La plataforma u otras personas hacen mal uso del context de alta dimensionalidad & Almacenamiento local, cifrado y autorización por niveles. \\
Socialización falsa & Las interacciones con IA ahogan las interacciones entre personas reales & Señalar qué es IA y qué es una persona real; optimizar la tasa de conexión entre personas reales. \\
Manipulación y adicción & El doble digital genera dependencia emocional para aumentar la retención & Objetivos transparentes y límites de uso saludable. \\
Emparejamiento erróneo & El emparejamiento de alta dimensionalidad produce una falsa sensación de intimidad & Confirmación gradual, retroalimentación de personas reales y avisos sobre límites. \\
\bottomrule
\end{tabularx}

\begin{warningbox}{Un mundo de bajo daño no es un resultado automático}
La IA reduce la fricción, pero también puede reducir las defensas. El producto debe incorporar desde el diseño los permisos, la transparencia, la confirmación por parte de personas reales y los límites de uso saludable; de lo contrario, el mundo de bajo daño puede convertirse en un mundo de alto daño más difícil de detectar.
\end{warningbox}

\subsection{Resumen de la sección}

Los riesgos y el valor de la socialización mediada por IA provienen de lo mismo: te entiende mejor y también actúa con más iniciativa. Solo al comprender esto se pueden diseñar productos que tengan capacidad de conexión y, a la vez, sentido de los límites.
\section{Modelo de datos del context: qué debe almacenar realmente la socialización con IA}

Para que el avatar se parezca al usuario, una ficha de perfil no basta. La socialización con IA necesita un context de varias capas: identidad estable, preferencias de largo plazo, estado reciente, estilo de expresión, objetivos sociales, grafo de relaciones, límites de privacidad e historial de retroalimentación. Cada capa tiene una vigencia temporal, una sensibilidad y una forma de autorización distintas, por lo que no pueden mezclarse en un único campo de «memoria».

\subsection{Capas del context}

\noindent\begin{longtable}{p{0.22\textwidth}p{0.34\textwidth}p{0.35\textwidth}}
\toprule
Capa & Contenido & Tratamiento en el producto \\
\midrule
Capa de identidad & Identidad básica, profesión, ciudad, profile público & Puede ser pública y editable; sensibilidad baja. \\
Capa de preferencias & Intereses, gusto estético, formas de expresión que agradan o desagradan & Se usa para las recomendaciones y para generar el tono. \\
Capa de relaciones & Amigos, seguidos, interacciones previas, vínculos débiles & Alto valor y también alta sensibilidad; requiere permisos. \\
Capa de estado & Emociones recientes, objetivos, actividades en curso & Muy sujeta al tiempo; requiere un mecanismo de caducidad. \\
Capa de expresión & Estilo de escritura, humor, opiniones, límites & Determina si el avatar se parece o no al usuario. \\
Capa de retroalimentación & Aprobación o rechazo del usuario a las acciones del avatar & Base para entrenar al avatar y ajustar la estrategia. \\
\bottomrule
\end{longtable}

\begin{importantbox}{Con el context no importa tener más, sino tenerlo mejor separado en capas}
Cada capa del context difiere en sensibilidad, vigencia temporal y uso. Si un producto de socialización con IA no las gestiona por separado, enfrentará al mismo tiempo riesgos de privacidad y una experiencia confusa.
\end{importantbox}

\subsection{Ciclo de vida del context}

El context debe tener un ciclo de vida: recopilación, interpretación, uso, retroalimentación, caducidad y eliminación. Por ejemplo, «hoy quiero conocer a gente que hace productos de IA» es un objetivo de corto plazo y no debe fijarse de forma permanente; «me desagrada el humor ofensivo» es una preferencia de largo plazo y puede conservarse por mucho tiempo; en cambio, «el contenido de cierta conversación privada» quizá solo pueda usarse de forma local y no debe entrar en el emparejamiento público.

\begin{warningbox}{El problema de la contaminación de la memoria}
Si el avatar toma una expresión momentánea como una preferencia de largo plazo, o una broma como un valor personal, se produce contaminación de la memoria. Los productos de socialización con IA deben permitir que el usuario corrija y elimine el context; de lo contrario, cuanto más se usen, menos se parecerán a él.
\end{warningbox}

\subsection{Resumen de la sección}

El context es el modelo de datos subyacente de la socialización con IA. Requiere separación en capas, autorización, caducidad y retroalimentación, no una simple acumulación de historiales de chat.
\section{Sistema de métricas: qué más observar además de la tasa de conexión entre personas reales}

Tristan menciona que la estrella polar es la tasa de conexión entre personas reales, y esa dirección es muy acertada. Sin embargo, un producto social no puede guiarse por una sola métrica. La tasa de conexión entre personas reales debe analizarse junto con la retención, la tasa de representación errónea, el grado de completitud del context, la tasa de revocación de privacidad, la proporción de interacciones IA/personas reales y la calidad de las relaciones a largo plazo.

\subsection{Tabla de métricas}

\noindent\begin{longtable}{p{0.24\textwidth}p{0.32\textwidth}p{0.35\textwidth}}
\toprule
Métrica & Qué mide & Riesgo de mal uso \\
\midrule
Tasa de conexión entre personas reales & Proporción de interacciones efectivas entre personas reales propiciadas por la IA & Puede inducir a molestar en exceso. \\
Proporción de actividad de personas reales & Porcentaje de las interacciones de la red que hacen personas reales y no la IA & Si es demasiado baja, la red se vuelve una IA que habla consigo misma. \\
Grado de completitud del context & Si el usuario proporciona suficiente context utilizable & Perseguirlo en exceso invade la privacidad. \\
Tasa de representación errónea & Proporción de salidas del doble digital que el usuario rechaza & Requiere un registro minucioso y una corrección rápida. \\
Retención & Si el usuario sigue regresando & Puede inflarse con una actividad aparente que no es real. \\
Calidad de las relaciones & Si las conexiones son duraderas, valiosas y de bajo daño & Es difícil de cuantificar, pero es la más cercana a la visión. \\
\bottomrule
\end{longtable}

\begin{knowledgebox}{Por qué las métricas deben analizarse en conjunto}
Fijarse solo en la tasa de conexión entre personas reales puede fomentar que se atraiga a la gente de forma brusca; fijarse solo en la retención puede fomentar la adicción; fijarse solo en el número de interacciones con la IA puede generar una prosperidad ficticia. La IA social necesita un conjunto de métricas que se equilibren entre sí.
\end{knowledgebox}

\subsection{Resumen de la sección}

La tasa de conexión entre personas reales es la estrella polar, pero no el único tablero de control. Un buen sistema de métricas para la IA social debe considerar a la vez la conexión, la privacidad, la representación errónea y la calidad de las relaciones.
\section{Estrategia de lanzamiento: de círculos reducidos al efecto de red}

La difusión inicial de Elys provino de círculos reducidos, en particular de personas atentas a los nuevos paradigmas de la IA. Este modo de arranque tiene ventajas: los primeros usuarios entienden los conceptos de context, doble digital y proactive, y están dispuestos a tolerar un producto todavía tosco. La desventaja es que un círculo reducido no equivale al mercado masivo; al salir de ese círculo, los problemas de privacidad, los malentendidos, el abuso y las interacciones de baja calidad se amplifican.

\subsection{Ruta de lanzamiento en tres etapas}

\noindent\begin{tabularx}{\textwidth}{p{0.22\textwidth}p{0.34\textwidth}X}
\toprule
Etapa & Objetivo & Riesgos principales \\
\midrule
Círculo semilla & Validar el paradigma y el volante de inercia central & Los usuarios saben demasiado de IA y no representan al público general. \\
Comunidades verticales & Validar la tasa de conexión entre personas reales y la calidad de las relaciones & Aumenta la presión sobre los límites de la comunidad, la privacidad y la gobernanza. \\
Mercado masivo & Ampliar el efecto de red y la monetización & Ruido de IA, abuso, confianza y riesgos regulatorios. \\
\bottomrule
\end{tabularx}

\begin{importantbox}{Salir del círculo inicial no es solo un problema de crecimiento}
Cuando una red social de IA sale de su círculo inicial, cambia la estructura de riesgos. Los primeros usuarios tratan al doble digital como un experimento; los usuarios del público general pueden tomarlo por la persona misma. El producto debe establecer mecanismos de permisos, señalización, apelación y gobernanza antes de salir del círculo inicial.
\end{importantbox}

\subsection{Resumen de la sección}

Un producto como Elys no puede limitarse a buscar un crecimiento rápido. La estrategia de lanzamiento de una red social de IA debe construirse a la par de los mecanismos de confianza; de lo contrario, el efecto de red crecerá junto con los riesgos.
\section{Conexión con EP139 y EP136}

EP139 aborda la historia técnica de los Agent y destaca el proactive agent y el universal digital agent; EP136 aborda el modelo como OS y destaca que el modelo se convierte en la capa de planificación de tareas; EP135, por su parte, traslada ambos juicios a los productos sociales: si la IA puede completar tareas por iniciativa propia, también puede ayudar por iniciativa propia a que las personas establezcan relaciones. Si el modelo es el OS, la red social podría convertirse en la aplicación con el context personal más sensible sobre ese OS.

\begin{knowledgebox}{Marco de lectura conjunta de los tres episodios}
EP139 ofrece la definición técnica de Agent, EP136 ofrece el panorama industrial del modelo como OS y EP135 ofrece un caso concreto de producto de IA social. En conjunto responden a una pregunta: cómo pasa el modelo de responder preguntas a actuar en nombre de las personas.
\end{knowledgebox}

\subsection{Resumen de la sección}

EP135 extiende los dos episodios anteriores hacia la conversión en producto para el consumidor final. No discute la capacidad del modelo en sí, sino cómo se reescribirán las redes sociales cuando el modelo pueda actuar por iniciativa propia.
\section{Términos clave: índice de palabras clave de este episodio}

\begin{longtable}{p{0.24\textwidth}p{0.32\textwidth}p{0.35\textwidth}}
\toprule
Término & Explicación en una frase & Papel en este episodio \\
\midrule
Elys & Producto social de IA de Ziran Xuanze (自然选择) & Promueve la conexión entre personas reales mediante alter egos digitales. \\
Eve & Producto de IA de compañía y juego & Encarna la línea de producto de «los seres humanos están demasiado solos». \\
Context & Estado de alta dimensión del usuario: experiencias, preferencias, valores, obras, relaciones & Activo central de la IA social. \\
Alter ego digital & Agente de IA que actúa en la red en representación del usuario & Convierte el context en acciones sociales. \\
Subjetividad & Estructura expresable y computable de «quién soy» del usuario & Base para que el alter ego se parezca al usuario. \\
Proactive & La IA hace cosas por iniciativa propia en nombre del usuario & Núcleo del cambio de paradigma de los productos de IA. \\
Tasa de conexión entre personas reales & Proporción de conexiones efectivas que se producen entre personas reales & Métrica estrella del norte de Elys. \\
Mundo de bajo daño & Entorno de conexión ideal, con menos fricción y menos daño social & Visión del producto. \\
\bottomrule
\end{longtable}

\subsection{Resumen de la sección}

Todas las palabras clave de este episodio giran en torno a una pregunta: cómo pasa la IA de ser una herramienta a ser un alter ego que enfrenta el mundo en nombre del usuario, sin dejar de estar al servicio de la conexión entre personas reales.
\section{Síntesis y ampliación}

\subsection{Conclusiones centrales}

\begin{enumerate}
\item Lo central de la socialización con IA no es que la IA converse con el usuario para hacerle compañía, sino que el doble digital, con su context, cree por iniciativa propia conexiones entre personas reales.
\item El context es el activo central de la socialización con IA: la inteligencia quizá llegue a estar al alcance de todos por igual; el context, no.
\item La subjetividad es la expresión en ingeniería de «ser uno mismo» y es la base para que el doble digital pueda actuar en representación del usuario.
\item La privacidad y el context forman una relación de intercambio: el producto debe hacer que el usuario perciba un beneficio antes de abandonarlo.
\item Elys no es un Tinder con IA; su objetivo es el emparejamiento por context de alta dimensión y la socialización proactive.
\item Las empresas de modelos no ganan por naturaleza la socialización con IA; la ventaja defensiva de las empresas de aplicaciones está en el context, las relaciones y el lugar que el producto ocupa en la mente del usuario.
\end{enumerate}

\subsection{Preguntas abiertas}

\begin{itemize}
\item ¿Estarán dispuestos los usuarios a entregar suficiente context a una empresa comercial?
\item Cuando el doble digital se expresa en nombre del usuario, ¿cómo se establecen los mecanismos de autorización, revocación y responsabilidad?
\item ¿La socialización con IA aumentará las conexiones entre personas reales o producirá más interacciones ficticias?
\item Entre actores como ChatGPT, Meta, Doubao (豆包) y Elys, ¿quién tiene más probabilidades de obtener una ventaja de context a largo plazo?
\item ¿Cómo puede un producto proactive evitar convertirse en acoso o salirse de control?
\end{itemize}

\subsection{Lecturas adicionales}

\begin{itemize}
\item EP139, Historia técnica de los Agent: para entender el contexto técnico de los proactive agent.
\item EP136, El modelo como OS: para entender cómo el modelo se convierte en la capa de planificación de tareas.
\item Material sobre interacción humano-computadora, redes sociales, sistemas de recomendación y cómputo con preservación de la privacidad: para entender las bases multidisciplinarias de las redes sociales con IA.
\item La película \textit{Her} y el debate sobre los productos de AI companion: para entender la diferencia entre la narrativa de la compañía y los productos reales.
\end{itemize}

\begin{importantbox}{Valoración final}
La verdadera dificultad de la socialización con IA no es lograr que la IA sepa conversar, sino lograr que, respetando la privacidad y la subjetividad, convierta el context de una persona en mejores conexiones con personas reales. Si lo consigue, la unidad básica de las redes sociales pasará de ser la «cuenta» a ser el «doble digital capaz de actuar».\end{importantbox}

\end{document}
